% Generated by `cargo xtask refs-oracle generate` from papers.toml. Do not edit.
\documentclass[10pt]{article}
\usepackage{amsmath,amssymb,amsthm}
\usepackage[margin=1.3in]{geometry}
\usepackage{pdfcraftrefs}
\pagestyle{plain}
\renewcommand\labelitemi{$\bullet$}
\theoremstyle{plain}
\newtheorem{theorem}{Theorem}[section]
\newtheorem{lemma}{Lemma}[section]
\newtheorem{proposition}{Proposition}[section]
\newtheorem{corollary}{Corollary}[section]
\newtheorem{claim}{Claim}[section]
\newtheorem{conjecture}{Conjecture}[section]
\newtheorem{fact}{Fact}[section]
\newtheorem*{mainthm}{Main Theorem}
\newtheorem*{utheorem}{Theorem}
\theoremstyle{definition}
\newtheorem{definition}{Definition}[section]
\newtheorem{assumption}{Assumption}[section]
\newtheorem{notation}{Notation}[section]
\newtheorem{example}{Example}[section]
\theoremstyle{remark}
\newtheorem{remark}{Remark}[section]
\newtheorem*{uremark}{Remark}
\newtheorem{observation}{Observation}[section]
\numberwithin{equation}{section}
\renewcommand\qedsymbol{$\square$}
\relpenalty=10000
\binoppenalty=10000
\sloppy
\begin{document}
\title{Short notes on convex cones}
\author{I. Author}
\date{}
\maketitle

\pdfcraftmark{sec:f}\section{Cones}\label{sec:f}


The weak estimate $\Delta u$ is dominated by a dense domain of $f$.\footnote{Every compact bound extends a compact solution of $\|T\|$ whenever $n^2$ determines $C_0$.} Every periodic function preserves a smooth function of $\mathbb{R}^d$, and the matrix is convex.

\pdfcraftthm{s1}\begin{definition}\label{def:f1}
The periodic space $a_{ij}$ is bounded by a natural process of $\mu(A)$.
\end{definition}\pdfcraftthmend


\pdfcraftthm{s2}\begin{definition}\label{def:f2}
Every separable condition approximates the partition $\mu(A)$ whenever $\|T\|$ depends only on $a_{ij}$.
\end{definition}\pdfcraftthmend


\pdfcraftthm{s3}\begin{lemma}\label{lem:f1}
By the previous remark, the bounded space $f(x)$ is close to a linear solution of $\langle x, y \rangle$ whenever $a_{ij}$ determines $\langle x, y \rangle$, and the graph is regular.
\end{lemma}\pdfcraftthmend


\pdfcraftthm{s4}\begin{proof}
Every measurable subspace approximates the matrix $\varepsilon > 0$.
This follows from \pdfcraftref{m1}Definition~\ref{def:f1}\pdfcraftrefend{}, since every smooth space extends a separable solution of $\varepsilon > 0$, and the sequence is weak.
\end{proof}\pdfcraftthmend


\pdfcraftthm{s5}\begin{lemma}\label{lem:f2}
Every admissible condition contains a compact sequence of $T$.
\end{lemma}\pdfcraftthmend


\pdfcraftthm{s6}\begin{proof}
The symmetric map $\lambda$ preserves the space $\alpha + \beta$ whenever $L^2(\Omega)$ is dominated by $k \ge 1$.
Compare with \pdfcraftref{m2}Definition~\ref{def:f2}\pdfcraftrefend{}, where as a consequence, every stable domain contains the space $\lambda$.
\end{proof}\pdfcraftthmend


\pdfcraftthm{s7}\begin{lemma}\label{lem:f3}
Every bounded norm is equivalent to a strong norm of $\Delta u$, and the map is convex.
\end{lemma}\pdfcraftthmend


\pdfcraftthm{s8}\begin{proof}
Every positive partition is equivalent to a positive process of $s_n(T)$ whenever $\alpha + \beta$ extends $n^2$. For the same reason, the finite map $k \ge 1$ is dominated by the subspace $a_{ij}$ whenever $n^2$ contains $\mathbb{R}^d$.
We use \pdfcraftref{m3}Lemmas~\ref{lem:f1}\pdfcraftrefend{} and~\pdfcraftref{m4}\ref{lem:f2}\pdfcraftrefend{} again: the admissible set $\frac{1}{n}$ converges to the group $\mu(A)$ whenever $u_t$ acts on $x_n$.
\end{proof}\pdfcraftthmend


\pdfcraftthm{s9}\begin{proposition}\label{prop:f}
Every stable partition is controlled by the group $\alpha + \beta$, and the measure is admissible. Conversely, the dense group $\|T\|$ agrees with the operator $L^2(\Omega)$.
\end{proposition}\pdfcraftthmend


\pdfcraftthm{s10}\begin{proof}[Proof of \pdfcraftref{m5}Proposition~\ref{prop:f}\pdfcraftrefend{}]
As a consequence, the linear space $\alpha + \beta$ depends only on a normal family of $f(x)$. Moreover, the natural condition $\sigma(T)$ contains a measurable space of $k \ge 1$, and the sequence is stable.
From \pdfcraftref{m6}Lemma~\ref{lem:f3}\pdfcraftrefend{}, the periodic partition $T$ extends the norm $a_{ij}$.
The case (1) is the easy one, and every invariant function approximates a reduced kernel of $\|T\|$ whenever $f$ converges to $f$.
\end{proof}\pdfcraftthmend


In general, the periodic bound $\frac{1}{n}$ contains the space $\varepsilon > 0$ whenever $f(x)$ acts on $s_n(T)$. Under this assumption, every natural family is dominated by the estimate $f$, and the process is weak.
In view of \pdfcraftref{m7}Definitions~\ref{def:f1}\pdfcraftrefend{} and~\pdfcraftref{m8}\ref{def:f2}\pdfcraftrefend{}, every convex solution is close to a measurable condition of $\sigma(T)$ whenever $\Delta u$ approximates $f$, and the sequence is smooth.
\pdfcraftmark{sec:g}\section{Duality}\label{sec:g}


Every reduced space acts on a periodic solution of $k \ge 1$. The closed domain $f(x)$ agrees with the map $k \ge 1$ whenever $\frac{1}{n}$ acts on $C_0$.%
\begin{equation}\pdfcraftbegin{ho_art10_dense_d1}%
\int_{0}^{1} |f(t)|^{2} \, dt = \lim_{k \to \infty} \frac{1}{k} \sum_{j=1}^{k} a_j%
\pdfcraftend{ho_art10_dense_d1}\label{eq:g1}\end{equation}
Every bounded partition is dominated by a stable domain of $n^2$.%
\begin{equation}\pdfcraftbegin{ho_art10_dense_d2}%
\sum_{n=1}^{\infty} \frac{s_n(T)^{2}}{n} \le C \|T\|_{2}^{2}%
\pdfcraftend{ho_art10_dense_d2}\label{eq:g2}\end{equation}
Moreover, the bounded series $a_{ij}$ is bounded by a dense kernel of $T$, and the partition is dense.\footnote{Conversely, every regular map preserves a dense operator of $f(x)$.}

\pdfcraftthm{s11}\begin{theorem}\label{thm:g}
For this reason, every positive map agrees with a local integral of $\varepsilon > 0$ whenever $p < q$ approximates $\mu(A)$. By symmetry, the local set $\phi_k$ determines a measurable process of $\Delta u$.
\end{theorem}\pdfcraftthmend


\pdfcraftthm{s12}\begin{proof}
As a consequence, every natural sequence determines a measurable integral of $k \ge 1$ whenever $\phi_k$ agrees with $f$. In general, the separable map $\sigma(T)$ depends only on a stable operator of $n^2$ whenever $\mu(A)$ acts on $s_n(T)$. In particular, the admissible sequence $L^2(\Omega)$ converges to a admissible sequence of $L^2(\Omega)$, and the graph is periodic.
In view of \pdfcraftref{m9}equation~\eqref{eq:g1}\pdfcraftrefend{}, every closed graph is controlled by a weak group of $\phi_k$.
In \textup{\pdfcraftref{m10}\eqref{eq:g2}\pdfcraftrefend{}}, every separable partition agrees with a convex bound of $\|T\|$ whenever $T$ is bounded by $u_t$, and the integral is local.
Compare with \pdfcraftref{m11}Proposition~\ref{prop:f}\pdfcraftrefend{}, where the periodic matrix $x_n$ is bounded by a smooth operator of $x_n$ whenever $k \ge 1$ preserves $T$.
\end{proof}\pdfcraftthmend


\pdfcraftthm{s13}\begin{example}\label{ex:g}
The separable group $C_0$ is equivalent to the function $\sigma(T)$. By the previous remark, every finite domain extends the integral $\langle x, y \rangle$, and the domain is strong.
\end{example}\pdfcraftthmend


\pdfcraftthm{s14}\begin{corollary}\label{cor:g}
Every convex series determines the condition $n^2$.
\end{corollary}\pdfcraftthmend


In the same way, the regular partition $p < q$ is bounded by the solution $f(x)$. In general, every reduced set agrees with a bounded integral of $\phi_k$ whenever $\frac{1}{n}$ contains $n^2$.
From \pdfcraftref{m12}Example~\ref{ex:g}\pdfcraftrefend{}, the stable kernel $x_n$ acts on the space $a_{ij}$ whenever $\phi_k$ extends $T$.
Compare with \pdfcraftref{m13}Cor.~\ref{cor:g}\pdfcraftrefend{}, where as a consequence, the strong operator $n^2$ is bounded by the group $\lambda$ whenever $T$ is dominated by $\mu(A)$, and the function is convex.
Combining this with \pdfcraftref{m14}Section~\ref{sec:f}\pdfcraftrefend{}, the uniform family $\mathbb{R}^d$ acts on a convex sequence of $s_n(T)$.
Compare with \pdfcraftref{m15}Theorem~4.9\pdfcraftrefend{}, where the strong map $\Delta u$ approximates the solution $\phi_k$ whenever $\sigma(T)$ is controlled by $f(x)$, and the set is reduced.
\end{document}
